\section{The same singular constants in the complete marked receiver}
We now prove the exact change of coordinates from RS27 to RD19,
including the mixed metric. The original reference period and its
distinct quartet are fixed throughout. Its ordered labels remain
$(\mathrm M,\mathrm N,\mathrm T,\mathrm E)$, with both signs at
each label. The varying quartic is RS1; its parameter $w=\delta^2$
does not replace that reference period. The marked frames
$K,X,Y,\Phi,\Psi,E,O$ are exactly FC26--32 and SF2--3.
In particular $Y=\Psi K$, $X=\Phi K$, $Y=T_AX$ and
$\Psi=T_A\Phi$. All their original phases and entries are retained.
The human polynomial provenance remains Alp\"oge's announcement
and Tao's account \cite{Tao}; the four labels and their unchanged
marking come from the ES source \cite{ES488}.

For a scalar $a$, define the complete cubic substitution matrix
\begin{equation}
 C_a=\begin{pmatrix}1&a&a^2&a^3\\0&1&2a&3a^2\\
 0&0&1&3a\\0&0&0&1\end{pmatrix},\qquad
 [p(x+a)]=C_a[p(x)],\quad C_aC_b=C_{a+b},\quad C_a^{-1}=C_{-a}.
 \tag{MR1}
\end{equation}
The coefficient of $x^i$ in $(x+a)^j$ is
$\binom ji a^{j-i}$, proving the first identity. Substitution
twice proves the composition and inverse without discarding
any original center. Retain $c_*=a_{\rm amp}/2-4$ and
$\kappa=c_*-1/2$. In RS1 the coordinate is $x=r-1/2$, so
the full $r$-coefficient matrix $Q_r$ of RD19 becomes
\begin{equation}
 \begin{split}
 Z&=C_{c_*}\Psi,\\
 Q_x&=\operatorname{diag}(C_{1/2},C_{1/2})Q_r
   =\begin{pmatrix}C_{c_*}Y&0\\C_{c_*}\Psi E&C_{c_*}\Psi O\end{pmatrix}
   =\operatorname{diag}(Z,Z)\begin{pmatrix}K&0\\E&O\end{pmatrix},\\
 Q_x^{-1}&=\begin{pmatrix}
 K^{-1}Z^{-1}&0\\-O^{-1}EK^{-1}Z^{-1}&O^{-1}Z^{-1}
 \end{pmatrix},\qquad
 \det Q_x=(\det\Psi)^2\det K\det O\ne0.
 \end{split}\tag{MR2}
\end{equation}
Indeed $C_{1/2}C_\kappa=C_{c_*}$ and $Y=\Psi K$.
Multiplication of the two displayed block matrices proves
both inverse products. The original odd frame $O$ is
invertible by SE6; all remaining factors are the original
FC isomorphisms. These facts prove the last assertion with
the actual determinant, rather than selecting a new frame.

Use the two coefficient receivers $\mathcal C_U,\mathcal C_W$
of RS27. The exact identities with the original marked maps are
\begin{equation}
 \mathcal J_U=\mathcal C_UQ_x,\qquad
 \mathcal J_W=\mathcal C_WQ_x,\qquad
 (T_A\oplus T_A)\mathcal J_U=\mathcal J_W.
 \tag{MR3}
\end{equation}
For example, $C_WC_{c_*}Y\alpha=Y\alpha$ since
$(C_{c_*}p)(S'-c_*)=p(S')$. Applying the original $T_A^{-1}$
gives $X\alpha$. The second component gives
$\Phi(E\alpha+O\beta)$ by the same substitutions. This proves
each component of MR3 and both directions of the correspondence.

Write $G_{U,x}$ for the complete RS24 source Gram in the
$x$-coefficient basis and $G_{W,x}$ for the corresponding
original target Gram. Reserve $G_U=\Phi^*\langle\ ,\ \rangle_U\Phi$
for its earlier marked meaning; these two matrices are related by
\begin{equation}
 \begin{gathered}
 G_U=Z^*G_{U,x}Z,\quad G_X=K^*G_UK,\qquad
 G=Z^*G_{W,x}Z,\quad G_Y=K^*GK,\\
 \mathbb H_U=Q_x^*\operatorname{diag}(G_{U,x},G_{U,x})Q_x
 =\begin{pmatrix}
 G_X+E^*G_UE&E^*G_UO\\O^*G_UE&O^*G_UO
 \end{pmatrix},\\
 \mathbb H_W=Q_x^*\operatorname{diag}(G_{W,x},G_{W,x})Q_x
 =\begin{pmatrix}
 G_Y+E^*GE&E^*GO\\O^*GE&O^*GO
 \end{pmatrix}.
 \end{gathered}\tag{MR4}
\end{equation}
The first line follows from $C_UZ=\Phi$, $C_UZK=X$ and
their target counterparts, proved in MR3. Multiplication
of the complete block matrices then proves every remaining
entry. In particular the mixed terms $E^*G_UO$ and $O^*G_UE$
are present; neither parity orthogonality nor a diagonal
marked metric was assumed. The full moment and Gamma formula
for $G_{U,x}$, its inverse and all phases is RS21--26.

Let $\Theta_w=M_{2t^3}/c$ be RS3 with its exact phase
$c=(1+i)/\sqrt2$, and put
\begin{equation}
 \Theta_{{\rm mark},w}=Q_x^{-1}\Theta_wQ_x,\qquad
 \Theta_{{\rm mark},w}^{-1}=Q_x^{-1}\Theta_w^{-1}Q_x.
 \tag{MR5}
\end{equation}
All maps in this identity act on the same eight signed
coordinates as RD19. Their physical original-source and
target forms are precisely $\mathcal J_U\Theta_{{\rm mark},w}
\mathcal J_U^{-1}$ and $\mathcal J_W\Theta_{{\rm mark},w}
\mathcal J_W^{-1}$. MR3 proves their intertwining by
$T_A\oplus T_A$ in both directions.

Here is the full singular-value argument including the metric.
For any invertible $Q$, positive $G_8$, and matrix $A$, set
$H=Q^*G_8Q$ and $A_Q=Q^{-1}AQ$. Direct inversion and
multiplication give
\begin{equation}
 H^{-1}A_Q^*HA_Q
       =Q^{-1}G_8^{-1}A^*G_8AQ.
 \tag{MR6}
\end{equation}
Both sides therefore have the same eigenvalues. Those eigenvalues
are the squared singular values in their respective positive
metrics, because conjugation by the positive square root of
the metric gives the usual adjoint product. Thus every singular
value is retained. Exterior norms are retained as well: the
Gram of wedges is the determinant of the individual pairings,
and an isometry preserves each of these pairings. With MR4,
this proves, for every $w>0$ and $0\le j\le8$,
\begin{equation}
 \left\|\bigwedge^j\Theta_{{\rm mark},w}^{-1}\right\|_{\mathbb H_U}
 =\left\|\bigwedge^j\Theta_w^{-1}\right\|_{
             \operatorname{diag}(G_{U,x},G_{U,x})}
 =\left\|\bigwedge^j\Theta_{U,w}^{-1}\right\|_{U\oplus U}.
 \tag{MR7}
\end{equation}
Consequently RS14--19 give the exact marked constants, with
the explicit original-source substitution RS24--25. The
target statement follows by the same proved identity with
$G_{W,x}$. No condition-number estimate replaces either
source or target constants in MR7.

\section{The complete nilpotent flag of the inverse pole}
Keep $g=\gamma^2>0$, $X_0,N,H_0$ of RS5, and let
\begin{equation}
 \begin{gathered}
 v_1=(g,0,1,0)^{\mathsf T},\quad
 v_2=(0,g,0,1)^{\mathsf T},\quad
 V=\operatorname{span}\{v_1,v_2\}=\operatorname{im}N=\ker N,\\
 \mathscr L=\lim_{w\downarrow0}w\Theta_w^{-1}
 =c\begin{pmatrix}0&H_0/(32g^2)\\(gI+N)/(32g^3)&0\end{pmatrix},\\
 \Theta_0=\frac2c\begin{pmatrix}0&0\\H_0&0\end{pmatrix}.
 \end{gathered}\tag{MR8}
\end{equation}
The first two columns of $N$ are $v_1,v_2$; its remaining
columns are $-gv_1,-gv_2$. They are independent, and $N^2=0$,
proving both descriptions of $V$. Its polynomials are exactly
$x^2+g$ and $x(x^2+g)$, with their original coefficients.
Since $X_0$ commutes with $N$ and is invertible,
$\operatorname{im}H_0=\ker H_0=V$ and $H_0N=NH_0=0$.
Multiplication of the unchanged matrices gives
\begin{equation}
 \mathscr L^2=\frac{c^2}{2^{10}g^4}
            \begin{pmatrix}H_0&0\\0&H_0\end{pmatrix},\qquad
 \mathscr L^3=\frac{c^3}{2^{15}g^6}
            \begin{pmatrix}0&0\\H_0&0\end{pmatrix}
       =-\frac{\Theta_0}{2^{16}g^6},\qquad
 \mathscr L^4=0.
 \tag{MR9}
\end{equation}
For the first product use $H_0(gI+N)=gH_0$; for the second
also use $H_0^2=0$. In comparing with $\Theta_0$ the factor
is $c^4/2^{16}g^6=-1/2^{16}g^6$, so the phase is essential.

The complete image and kernel flags, in even then odd order, are
\begin{equation}
 \begin{array}{c|c|c|c}
 j&\operatorname{im}\mathscr L^j&\ker\mathscr L^j&
     \operatorname{rank}\mathscr L^j\\\hline
 0&\C^4\oplus\C^4&0&8\\
 1&V\oplus\C^4&0\oplus V&6\\
 2&V\oplus V&V\oplus V&4\\
 3&0\oplus V&V\oplus\C^4&2\\
 4&0&\C^4\oplus\C^4&0
 \end{array}
 \tag{MR10}
\end{equation}
Indeed $(gI+N)^{-1}=g^{-1}I-g^{-2}N$ since $N^2=0$,
whereas $H_0$ has precisely image and kernel $V$. This proves
the first row of nontrivial subspaces directly from MR8;
MR9 proves the others. In particular
$\ker\Theta_0=V\oplus\C^4$ and
$\operatorname{im}\Theta_0=0\oplus V$.

There are two explicit complete chains. Let $e_0,e_1$
denote the first two cubic coefficient vectors, and set
$u_j=(e_j,0)$, for $j=0,1$. Without rescaling any vector they are
\begin{equation}
 \begin{split}
 u_j&=(e_j,0),\\
 \mathscr L u_j&=\left(0,\frac{c(gI+N)e_j}{32g^3}\right),\\
 \mathscr L^2u_j&=\left(\frac{c^2H_0e_j}{1024g^4},0\right),\\
 \mathscr L^3u_j&=\left(0,\frac{c^3H_0e_j}{32768g^6}\right),\\
 H_0e_0&=-2v_2,\qquad H_0e_1=2gv_1.
 \end{split}\tag{MR11}
\end{equation}
Both last vectors are independent. In any linear relation
between these eight vectors, applying $\mathscr L^3$ first
forces the coefficients of $u_0,u_1$ to vanish. Applying
$\mathscr L^2$ then forces the coefficients of
$\mathscr L u_0,\mathscr L u_1$ to vanish; applying
$\mathscr L$ and then using the independence of the last
vectors eliminates the remaining coefficients. Thus the
ordered columns
$(\mathscr L^3u_0,\mathscr L^2u_0,\mathscr L u_0,u_0,
  \mathscr L^3u_1,\mathscr L^2u_1,\mathscr L u_1,u_1)$
are a basis and exhibit exactly two blocks of length four.

Every vector and subspace in this flag returns to the original
marked and physical spaces by the explicit isomorphisms
\begin{equation}
 \mathscr L_{\rm mark}=Q_x^{-1}\mathscr LQ_x,
 \quad \operatorname{im}\mathscr L_{\rm mark}^j
       =Q_x^{-1}\operatorname{im}\mathscr L^j,
 \quad \ker\mathscr L_{\rm mark}^j
       =Q_x^{-1}\ker\mathscr L^j,
 \quad \mathcal J_UQ_x^{-1}=\mathcal C_U.
 \tag{MR12}
\end{equation}
The corresponding target map is $\mathcal J_WQ_x^{-1}
=\mathcal C_W$. Applying $T_A\oplus T_A$ to the source
chain gives the target chain term by term. MR2 supplies the
entire inverse, so no marked state or phase is omitted.
For any chain column matrix $B_j$, its actual source Gram
is exactly $B_j^*\operatorname{diag}(G_{U,x},G_{U,x})B_j$,
which equals $(Q_x^{-1}B_j)^*\mathbb H_U(Q_x^{-1}B_j)$
by MR4. This also proves the metric statement for every flag.

\section{Exact finite-parameter inverse bounds with the original metric}
The singular asymptotes admit a fully evaluated error term.
For every $w>0$, the matrices RS4 satisfy
\begin{equation}
 \Theta_w^{-1}=\frac{\mathscr L}{w}
       +c\begin{pmatrix}0&A_w\\B_w&0\end{pmatrix},
 \tag{MR13}
\end{equation}
where every entry of the remainder is
\begin{equation}
 A_w=\begin{pmatrix}
 0&-\frac1{16g}&0&-\frac{2g+w}{16g}\\
 \frac1{2(g+w)^2}&0&-\frac1{16g}&0\\
 0&0&0&\frac1{16g}\\
 \frac{3g+w}{16g^2(g+w)^2}&0&0&0
 \end{pmatrix},\qquad
 B_w=\begin{pmatrix}
 -\frac{3g+w}{16g^2(g+w)^2}&0&0&0\\
 0&-\frac{3g+w}{16g^2(g+w)^2}&0&0\\
 -\frac{2g+w}{32g^3(g+w)^2}&0&0&0\\
 0&-\frac{2g+w}{32g^3(g+w)^2}&0&0
 \end{pmatrix}.
 \tag{MR14}
\end{equation}
For a direct verification put $a=g-w$, $p=(g+w)^2$.
The exact inverse of RS2 is
\begin{equation}
 X_w^{-1}=\begin{pmatrix}
 0&1&0&0\\-2a/p&0&1&0\\0&0&0&1\\-1/p&0&0&0
 \end{pmatrix}.
 \tag{MR15}
\end{equation}
Multiplying both products with $X_w$ proves this inverse.
Subtracting $H_0/(32g^2w)$ from
$(X_w+aX_w^{-1})/(16gw)$ gives exactly $A_w$.
Subtracting $(gI+N)/(32g^3w)$ from
$-X_w^{-2}/(32gw)$ gives exactly $B_w$.
These are entrywise calculations and prove MR13--14, including
the factors of $g$, $w$ and $c$. Every entry extends continuously
to $w=0$ for the unchanged $g>0$.

For the original source set $G_0=G_{U,x}$ with its complete
RS24--25 evaluation, and define the following explicit finite
quantities:
\begin{equation}
 \begin{split}
 e(w)^2&=\operatorname{tr}(G_0^{-1}A_w^*G_0A_w)
            +\operatorname{tr}(G_0^{-1}B_w^*G_0B_w),\\
 b(w)^2&=4\operatorname{tr}(G_0^{-1}(H_w^2)^*G_0H_w^2)
            +4\operatorname{tr}(G_0^{-1}H_w^*G_0H_w),\\
 p_j(w)&=\max\{0,c_j-we(w)\},\qquad
 q_j(w)=c_j+we(w),\quad 1\le j\le6 .
 \end{split}\tag{MR16}
\end{equation}
The $c_j$ are the six explicitly evaluated RS14 constants
for this $G_0$. All entries in MR16 are displayed in
RS2, RS24--25 and MR14. The two traces in each line are
sums of squares in the original positive metric. Hence $e(w)$
is the full Hilbert--Schmidt norm of the remainder and $b(w)$
is that of $\Theta_w$. Both are finite, and $b(w)>0$ for
every $w>0$. The phase has modulus one and remains present
in MR13 even though it cancels from these norms.

The variational characterization of singular values gives
$|\sigma_j(A)-\sigma_j(B)|\le\|A-B\|$: apply the triangle
inequality on each subspace in the minimum-maximum formula,
and then reverse $A,B$. Apply this to $w\Theta_w^{-1}$ and
$\mathscr L$. Its first six singular values are $c_j$, and
its last two are zero. The operator norm is bounded by the
Hilbert--Schmidt norm, so MR13 proves
\begin{equation}
 p_j(w)\le w\sigma_j(\Theta_{{\rm mark},w}^{-1})\le q_j(w)
 \quad(1\le j\le6),\qquad
 b(w)^{-1}\le\sigma_j(\Theta_{{\rm mark},w}^{-1})\le e(w)
 \quad(j=7,8).
 \tag{MR17}
\end{equation}
For the lower bound on the last two, the smallest singular
value of an inverse is the reciprocal of the largest singular
value of the original. This follows by inverting orthonormal
singular bases. The latter largest value is at most $b(w)$.
MR7 transfers all these inequalities to the actual marked
metric without any comparison loss.

Multiplying the decreasing singular values yields, for every
$w>0$, the complete inverse exterior bounds
\begin{equation}
 \begin{aligned}
 \frac{\prod_{j=1}^r p_j(w)}{w^r}
 &\le\|\bigwedge^r\Theta_{{\rm mark},w}^{-1}\|_{\mathbb H_U}
 \le\frac{\prod_{j=1}^r q_j(w)}{w^r},&&1\le r\le6,\\
 \frac{\prod_{j=1}^6p_j(w)}{w^6b(w)^{r-6}}
 &\le\|\bigwedge^r\Theta_{{\rm mark},w}^{-1}\|_{\mathbb H_U}
 \le\frac{\prod_{j=1}^6q_j(w)e(w)^{r-6}}{w^6},&&r=7,8,\\
 \|\bigwedge^8\Theta_{{\rm mark},w}^{-1}\|_{\mathbb H_U}
 &=\frac1{2^{32}g^6w^6(g+w)^6}.
 \end{aligned}\tag{MR18}
\end{equation}
Rank zero has norm one. Continuity of MR14 at zero implies
$we(w)\to0$; all $c_j$ are positive by RS8 and RS13.
Thus the lower bounds on the first six factors become positive
near zero and approach their exact constants. RS19 supplies
the two remaining exact leading constants. Equations MR7 and
MR18 are the full original-metric conclusion for the completed
trace-to-residue operation, at finite parameters as well as in
the limit.

The loss occurs in the explicit multiplication $2t^3/c$, whose
critical rank is two and whose inverse has the chains MR11.
The original conductor itself is the fixed isomorphism FC23--24.
MR3, MR5 and MR12 prove the exact relation to that isomorphism;
they do not replace its moment $\mu_v$ or change its order.
